\chapter*{Preface}
% \addcontentsline{chapter}{Preface}
I seek to present here what I take to be the main themes and problems of the politics of Marxism, which might also be called the Marxist approach to politics, or Marxist politics. I have attempted to do this by way of a 'theorization' of material drawn mainly from the writings of Marx, Engels, and Lenin, with occasional references to other major figures of Marxism; and I have also drawn on the actual experience of the politics of Marxism—or of what has been claimed to be such—in this century.

I am very conscious that I have not dealt with many things which would have to be included in a book that made any claim to being a comprehensive work on Marxism and politics. I make no such claim and did not set out to write such a book. But I think that the book which I have written may help the interested reader to see more clearly what are the distinctive features of Marxist politics, and what are its problems.

Whatever view may be taken of my text and its arguments, there cannot at least be much doubt that the questions discussed here have been absolutely central to the politics of the twentieth century, and that they are likely to dominate what remains of it.

\noindent{}April 1976 \hfill R.M.

% \begin{tabularx}{1\textwidth}{ 
%   >{\raggedright\arraybackslash}X 
%   >{\centering\arraybackslash}X }
%     April 1976 &  R.M. \\
% \end{tabularx}

\input{input/commands for pages end}